\subsection{Ideals of a valuation ring}

DEFINITION 2. Let G be an ordered set. A subset of G is called major if the relations \(x \in \mathbf{M}\) and \(y \geqslant x\) imply \(y \in \mathbf{M}\) .

Let K be a field, v a valuation on K, A the ring of v and G the order group of v. For every major subset \(M \subset G\) , let \(\mathfrak{a}(M)\) be the set of \(x \in K\) such that \(v(x) \in M \cup \{+\infty\}\) . Clearly \(\mathfrak{a}(M)\) is a sub-A-module of K.

PROPOSITION 7. The mapping \(\mathbf{M} \mapsto \mathfrak{a}(\mathbf{M})\) is an increasing bijection of the set of major subsets of \(\mathbf{G}\) onto the set of sub-A-modules of \(\mathbf{K}\) .

Let \(\mathfrak{b}\) be a sub-A-module of \(\mathbf{K}\) . The set of \(v(x)\) for \(x \in \mathfrak{b} - (0)\) is a major subset \(\mathbf{M}(\mathfrak{b})\) of \(\mathbf{G}\) . Proposition 7 will be shown if the following equations are proved:

\begin{enumerate}
\def\labelenumi{(\arabic{enumi})}
\setcounter{enumi}{1}
\item
  \(\mathbf{M}(\mathfrak{a}(\mathbf{N})) = \mathbf{N}\) for every major subset \(\mathbf{N}\) of \(\mathbf{G}\) ;
\item
  \(\mathfrak{a}(\mathbf{M}(\mathfrak{b})) = \mathfrak{b}\) for every sub-A-module \(\mathfrak{b}\) of \(\mathbf{K}\) .
\end{enumerate}

Formula (2) is easy, since, for all \(m \in \mathbf{N}\) , there exists \(x \in \mathbf{K}\) such that \(v(x) = m\) . Then obviously \(\mathfrak{b} \subset \mathfrak{a}(\mathbf{M}(\mathfrak{b}))\) ; conversely, let \(x \in \mathfrak{a}(\mathbf{M}(\mathfrak{b}))\) and suppose \(x \neq 0\) ;

then \(v(x) \in \mathbf{M}(\mathfrak{b})\) and therefore there exists \(y \in \mathfrak{b}\) such that \(v(x) = v(y)\) ; whence \(x = uy\) where \(v(u) = 0\) , which proves that \(x \in \mathrm{Ay} \subset \mathfrak{b}\) and completes the proof.

COROLLARY. Let \(G_{+}\) be the set of positive elements in G. The mapping \(\mathbf{M} \mapsto \mathfrak{a}(\mathbf{M})\) is a bijection of the set of major subsets of \(G_{+}\) onto the set of ideals of A.

As \(\mathbf{A} = \mathfrak{a}(\mathbf{G}_{+}),\mathfrak{a}(\mathbf{M})\subset \mathbf{A}\) is equivalent to \(\mathbf{M}\subset \mathbf{G}_{+}\) .

For example the maximal ideal \(m(A)\) is equal to \(a(S)\) , where S denotes the set of strictly positive elements of G.
% semantic-fragments: legacy-input-seam
\relax{}
